%%%% CAPÍTULO 1 - INTRODUÇÃO
\chapter{Introdução}\label{cap:introducao}

Enquanto o ensino tradicional foca na leitura e escrita de 
textos, os cursos de programação introduzem uma nova de forma de 
pensamento: os algoritmos. Estes são sequências lógicas 
interpretadas por humanos e máquinas. Essa disciplina, 
presente em cursos de áreas tecnológicas, é ``considerada 
desafiadora para alunos e professores"~ \cite{raabe2005ambiente}, 
pois exige o desenvolvimento de uma lógica formal.

Ainda segundo \cite{raabe2005ambiente}, as dificuldades no 
aprendizado de programação podem ser categorizadas em três grupos 
distintos, que vão além da simples aptidão individual: Problemas 
de natureza didática (relacionados aos métodos de ensino), 
cognitiva (relacionados aos desafios de raciocínio lógico e  
abstração) e afetiva~(relacionados à motivação, ansiedade e 
confiança do aluno). Essa categorização sistêmica ajuda a 
entender que o desafio requer iniciativas diferentes.  

No contexto educacional, a avaliação é entendida como uma etapa 
do processo de ensino e aprendizagem, o que a torna um tema 
permanente de reflexão e estudo entre especialistas da área, que 
buscam constantemente refinar suas práticas para alcançar 
melhores resultados \cite{lima2022avaliaccao}.

A avaliação no curso técnico configura-se como um processo 
bifronte, pois contempla tanto a avaliação da aprendizagem do 
estudante quanto a avaliação do sistema educacional como um todo. 
No que se refere à aprendizagem, o projeto pedagógico estabelece 
como objetivo central a progressão do discente para o 
alcance do perfil profissional almejado. Nesse sentido, o Art. 
34º da Resolução CNE/CEB nº 6, de 20 de setembro de 2012, que 
rege o curso, determina que esse processo ``é contínuo e 
cumulativo, com prevalência dos aspectos qualitativos 
sobre os quantitativos, bem como dos resultados ao longo do 
processo sobre os de eventuais provas finais" ~\cite{IFAM2020}.

Paralelamente, a avaliação do sistema, referente ao rendimento 
acadêmico, deve ser realizada para cada disciplina 
individualmente, considerando de forma simultânea e obrigatória a 
frequência do aluno e seu aproveitamento na aquisição de 
competências e conhecimentos \cite{IFAM2020}.

O ensino de programação é uma atividade multidisciplinar que 
demanda uma práxis pedagógica apurada --- ou seja, uma constante 
reflexão que integre teoria e prática em sala de aula. Uma 
estratégia comum é utilizar exemplos do cotidiano para estimular 
a resolução de problemas por meio de algoritmos. Contudo, 
avaliar a qualidade desses algoritmos exige ir além da 
verificação funcional. 

Esses algoritmos, definidos como procedimentos computacionais que 
transformam entradas em saídas, geram registros conhecidos como 
logs. Mas como esses logs podem ser utilizados para transformar o 
ensino e a aprendizagem? Este texto explora a interseção entre 
tecnologia e educação, mostrando como a análise de logs pode 
oferecer perspectivas valiosas para personalizar o ensino e 
identificar dificuldades dos alunos.

\begin{figure}[H]
\centering
\begin{tikzpicture}[node distance=2cm, auto, 
    block/.style={rectangle, draw, text width=3cm, align=center, minimum 
    height=1cm, rounded corners},
    arrow/.style={->, >=stealth, thick}]

    % Nós
    \node[block] (input) {INPUT \\ (Conteúdo ensinado)};
    \node[block, right=of input] (log) {LOG \\ (Rastro do processo)};
    \node[block, right=of log] (output) {OUTPUT \\ (Código final)};
    \node[block, below=of log, fill=gray!10] (analise) {Análise \\ pedagógica};

    % Setas
    \draw[arrow] (input) -- (log);
    \draw[arrow] (log) -- (output);
    \draw[arrow] (log) -- (analise);

\end{tikzpicture}
\caption{A posição dos logs na relação entre ensino (input) e produção discente 
(output)}
\label{fig:logs-intersecao}
\end{figure}

Esses registros, ou logs, são gerados continuamente durante as 
interações entre humanos e máquinas. Conforme trazido por 
\cite{clemente2008arquitetura}, ``log é definido como um 
conjunto de registros com marcação temporal, que suporta apenas 
inserção, e que representa eventos que aconteceram em um 
computador ou equipamento de rede". 
Além disso, segundo \cite{gu2022logging}, ``um log geralmente 
contém um grande número de entradas, também conhecidas como 
mensagens de registro". Em geral, o nível do log, seu conteúdo e 
o local adequado para declará-lo são decisões tomadas pelos 
desenvolvedores durante a criação do software.

O processo de ensino pode ser compreendido analogamente a uma 
tecnologia que pode ser atualizada, adaptada ou aplicada em 
diferentes contextos. No entanto, para estimar o tempo de 
produção de um código, métodos tradicionais de marcação 
de tempo são insuficientes, sendo necessárias métricas 
adicionais. Nesse sentido, o armazenamento de logs permite 
coletar informações detalhadas sobre eventos em dispositivos. 
Como destacado por \cite{nguyen_analyzing_2023}, 
``esse recurso de previsão possibilita que instrutores 
identifiquem alunos com dificuldades desde o início e ofereçam 
intervenções direcionadas para apoiar sua jornada de 
aprendizado". Ademais, a análise de logs pode revelar padrões 
de comportamento e envio, fornecendo, assim,  dados valiosos não 
apenas para a intervenção individual, mas para a melhoria 
contínua  e baseada em evidências do processo educacional como um 
todo.

A coleta automatizada de logs de programação para análise 
empírica é uma metodologia consolidada em estudos de Engenharia 
de Software. O projeto Besouro, por exemplo, foi concebido como 
um instrumento para capturar dados brutos do processo de TDD, com 
o objetivo de `comparar definições existentes, avaliá-las com 
relação à percepção dos usuários e explorar informações do 
código para análise posterior' \cite{becker2015besouro}. 
Inspirado por essa abordagem, o presente trabalho propõe a 
extensão LOGADO, que adapta esse princípio de instrumentação para 
o contexto educacional. Nosso plugin coleta logs de eventos de 
digitação (palavras reservadas, timestamps, localização no 
código). O objetivo principal é registrar os rastros da 
construção do código para, a partir da análise desses registros, 
encontrar formas de auxiliar o ensino de programação, oferecendo 
uma base de dados rica para pesquisas em ensino de computação.

\begin{figure}[H]
\centering
\begin{tikzpicture}[
    node distance=1.5cm,
    estudante/.style={rectangle, draw, minimum width=1.5cm, minimum height=1.8cm, align=center, rounded corners=0.3cm, fill=gray!5},
    plugin/.style={rectangle, draw, minimum width=2cm, minimum height=0.8cm, align=center, fill=white, rounded corners=0.05cm, thin},
    database/.style={cylinder, draw, shape aspect=0.4, minimum width=1.5cm, minimum height=2.2cm, align=center, cylinder uses custom fill, cylinder end fill=gray!20, cylinder body fill=white},
    arrow/.style={->, >=stealth, thick}
]

\node[estudante] (estA) {\faUser\\ Estudante A};
\node[estudante, right=of estA] (estB) {\faUser\\ Estudante B};
\node[estudante, right=of estB] (estN) {\faUser\\ Estudante N};

\node[plugin, below=0.3cm of estA] (pluginA) {LOGADO};
\node[plugin, below=0.3cm of estB] (pluginB) {LOGADO};
\node[plugin, below=0.3cm of estN] (pluginN) {LOGADO};

\node[database, below=1.8cm of estB] (db) {Dataset\\Agregado};

\node[rectangle, draw, below=1.2cm of db, minimum width=3cm, minimum height=0.8cm, align=center, fill=blue!10] (analise) {Análise Exploratória};

\draw[arrow] (pluginA.south) -- ++(0,-0.5) -| (db.west);
\draw[arrow] (pluginB.south) -- (db.north);
\draw[arrow] (pluginN.south) -- ++(0,-0.5) -| (db.east);
\draw[arrow] (db.south) -- (analise.north);

\node[below=0.5cm of pluginA, font=\small, xshift=-0.6cm] {logs (JSON)};
\node[below=0.15cm of pluginB, font=\small] {logs                  (JSON)};
\node[below=0.5cm of pluginN, font=\small, xshift=0.6cm] {logs (JSON)};

\end{tikzpicture}
\caption{Arquitetura de coleta e agregação de logs pelo plugin LOGADO.Fonte: Elaborado pelo autor.}
\label{fig:arquitetura-coleta}
\end{figure}

A necessidade de analisar o processo de aprendizagem, e não apenas seu produto final, tem sido reconhecida em diferentes domínios da computação. No ensino do paradigma orientado a objetos, por exemplo, \cite{felisbino2018supporting} desenvolveram uma ferramenta com um gerador de logs como plugin para o Microsoft Visio, justamente para capturar as ações dos alunos durante a construção de diagramas de classes e identificar suas dificuldades conceituais. Este trabalho se insere nessa mesma linha de investigação, transladando o princípio da coleta de logs processuais para o domínio do ensino de programação introdutória e de algoritmos.

Diferentemente da análise estática tradicional, que examina o código já pronto, a abordagem proposta registra o processo de construção do código em tempo real. É nesse contexto que métricas baseadas em rastros de construção, aplicadas com um viés pedagógico, se tornam relevantes.

Para além das métricas tradicionais de complexidade algorítmica (como tempo de execução e uso de memória), o docente pode apoiar-se em métricas que avaliam o processo de escrita do código, tais como a sequência de declaração de variáveis, o uso temporal de palavras reservadas e a alternância entre estruturas de controle. Com o suporte do plugin LOGADO, o docente passa a dispor de rastros detalhados do processo de escrita dos estudantes, algo que não seria viável registrar manualmente em um laboratório de informática com a mesma precisão. Esses rastros permitem que o docente identifique padrões de dificuldade e ofereça intervenções mais direcionadas, sem depender exclusivamente da observação presencial imediata.

O potencial inovador dessa abordagem reside na possibilidade de sistematizar a análise: os dados gerados por essas métricas, coletados em larga escala dos códigos dos estudantes, podem ser agregados em um dataset. Este, por sua vez, permite comparações objetivas não apenas entre indivíduos, mas entre turmas, períodos letivos e diferentes metodologias de ensino para avaliação em larga escala.
\section{Contexto, motivação e justificativa}\label{sec:consideracoesIniciais}

Diferentemente da avaliação quantitativa, que prioriza resultados finais, a avaliação qualitativa é fundamental para aferir o processo de aprendizado. Em cursos de algoritmos, esse acompanhamento processual pode ser significativamente ampliado com o uso da telemetria. Conforme proposto por \cite{kennedy2015towards}, essa abordagem é complementar aos mecanismos tradicionais de devolutiva, como notas em laboratórios, tarefas e exames. A telemetria não os substitui, mas oferece uma camada adicional de dados contínuos e objetivos sobre o comportamento do estudante, enriquecendo a análise do docente.

O valor desses dados é claro: ``saber o que acontece dentro de um sistema enquanto ele está em execução é benéfico para fins de depuração, manutenção e análise" \cite{Gatev2021}. Um log é um registro com marcação temporal que captura eventos de um sistema, formando um rico histórico de execução \cite{gu2022logging}.Nosso trabalho armazena logs dos discentes enquanto o código é produzido, salvando esses registros em arquivos externos ao código-fonte. Dessa forma, podemos observar suas estratégias sem comprometer o ambiente de desenvolvimento.

Embora a Mineração de Dados Educacionais tradicional pressuponha bases de dados estruturadas, a coleta por meio de logs nativos de programação, e não através de acessos em LMS, permanece pouco explorada na literatura. Uma das poucas exceções é o trabalho de \cite{kennedy2015towards}, que aborda a captura de telemetria em tutoriais de programação.

A telemetria é amplamente utilizada em ferramentas de desenvolvimento de software, mas sua aplicação no ensino de programação ainda é um campo em aberto. Ferramentas como o LOGADO contribuem para identificar dificuldades de estudantes a partir de dados quantitativos (timestamps) e qualitativos (palavras reservadas, linha e coluna do código). O estudante, como autor do texto performativo e instrucional, produz o código; o professor, como curador, analisa os logs para compreender o processo do aluno. O código-fonte segue como protagonista, mas os logs ganham destaque como deuteragonistas: pouco explorados até agora, esses registros revelam aspectos da codificação que métricas tradicionais deixam passar, possibilitando intervenções mais direcionadas e eficazes.

Conforme destacado por \cite{qian2017students}, ``esforços para aprimorar o ensino em ciência da computação estão em andamento, e os docentes enfrentam desafios significativos em disciplinas introdutórias de programação, visando facilitar a aprendizagem dos estudantes". Esse diagnóstico, embora formulado a partir da realidade universitária, encontra eco nos cursos técnicos em informática ofertados pela EPCT. É diante desse cenário que a coleta e análise de logs de programação se apresenta como uma via promissora: ao registrar timestamps, palavras reservadas e posições no código, é possível acompanhar o processo do estudante e não apenas o resultado final. Pesquisas como essa mostram que universidades e institutos federais, apesar de contextos distintos, compartilham o mesmo desafio e podem atuar juntos na busca por soluções, reforçando que a pesquisa produzida na EPCT dialoga com a produção acadêmica das universidades. Apesar dos esforços institucionais para mitigar esses problemas, disciplinas como Introdução a Algoritmos, Programação Básica, Estrutura de Dados e Programação Orientada a Objetos são frequentemente associadas a altas taxas de dificuldade e evasão. Além dos obstáculos cognitivos, os estudantes frequentemente lidam com questões emocionais, sentindo-se incapazes de dominar os conceitos e desenvolvendo, consequentemente, a crença de que são inaptos para a área de informática como um todo.

A coleta e análise de métricas, especificamente logs de atividade como timestamps, palavras reservadas e posições no código, em ambientes de desenvolvimento integrado (IDE) possibilitam avaliações mais robustas do processo de aprendizagem. A opção por um editor amplamente adotado no setor profissional, o Visual Studio Code (VS Code), que lidera o uso com 75,9\% dos desenvolvedores segundo o levantamento mais recente da indústria \cite{StackOverflowSurvey2025}, busca uniformizar as ferramentas utilizadas no contexto acadêmico, reduzindo a distância entre o que se ensina em disciplinas distintas e as exigências do mercado. A solução proposta consiste, portanto, em um plugin de operação transparente para essa plataforma, executando-se em segundo plano sem comprometer a experiência dos usuários finais (discentes e docentes) durante suas atividades de desenvolvimento.

\section{Questões de pesquisa}\label{sec:questoes-de-pesquisas}
\begin{enumerate}
    \item QP1: Quais são as principais aplicações do uso de logs no ensino de programação?
    \item QP2: Quais ferramentas, métodos e métricas são mais utilizados para capturar, armazenar e analisar logs em ambientes de ensino de programação?
\end{enumerate}

Para responder a essas questões, foi conduzido um Mapeamento Sistemático da 
Literatura (MSL), detalhado no Apêndice ~\ref{chap:msl}.
%msl-questoes
\section{Objetivos}
\label{sec:objetivos}

Este trabalho apresenta um plugin para coleta e análise de logs durante o processo de ensino-aprendizagem de programação. A ferramenta, já implementada e testada, fornece dados detalhados sobre os rastros de construção dos estudantes. Os logs coletados permitiram identificar dificuldades, gerar gráficos e observar características do processo de codificação, oferecendo subsídios para a personalização do ensino e para o aprimoramento das práticas pedagógicas. Os dados e a metodologia aqui apresentados também podem servir de referência para novos projetos na área de telemetria educacional.

\subsection{Objetivos Gerais}
\label{subsec:objetivos-gerais}

Esta pesquisa tem como objetivos gerais:

\begin{enumerate}
    \item Compreender, por meio da análise de logs acadêmicos, como estudantes de programação constroem suas soluções ao longo do tempo, identificando não apenas o que produzem, mas como produzem e quais estratégias, hesitações, correções e escolhas estilísticas emergem durante o processo de codificação.
    
    \item Avaliar a viabilidade e os limites do uso de telemetria no contexto de cursos técnicos de programação, investigando que tipo de inferências pedagógicas podem (e não podem) ser extraídas a partir dos rastros digitais coletados pelo plugin LOGADO.
    
    \item Oferecer à comunidade acadêmica um dataset de logs educacionais, coletado em condições reais de ensino, que possa servir de base para novas investigações sobre dificuldades de aprendizagem, padrões de erro e estratégias de resolução de problemas.
\end{enumerate}

\subsection{Tarefas Metodológicas}
\label{subsec:tarefas-metodologicas}

Para atingir os objetivos propostos, o trabalho se organiza nas seguintes tarefas metodológicas:

\begin{enumerate}
    \item Realizar um Mapeamento Sistemático da Literatura para identificar o estado da arte sobre o uso de logs no ensino de programação, mapeando ferramentas, métricas, aplicações e lacunas existentes.
    
    \item Projetar e implementar o plugin LOGADO para o Visual Studio Code, especificando quais eventos serão capturados (teclas, compilações, salvamentos, execuções) e justificando cada escolha com base nos objetivos de análise.
    
    \item Coletar logs em uma disciplina de programação do Instituto Federal do Amazonas, garantindo os aspectos éticos e de privacidade exigidos pela legislação (LGPD) e pelos comitês de ética.
    
    \item Processar e estruturar os dados brutos em um dataset organizado, documentando as decisões de limpeza, anonimização e agregação adotadas.
    
    \item Analisar os logs coletados sob duas perspectivas complementares: (a) uma análise descritiva, identificando padrões gerais de comportamento; (b) uma análise correlacional, explorando relações entre padrões de codificação e medidas de desempenho.
    
    \item Discutir as implicações pedagógicas dos achados, apontando tanto o potencial quanto as limitações da abordagem proposta para a prática docente em cursos de programação.
\end{enumerate}

\section{Contribuições esperadas}
\label{sec:contribuicoes}

Este trabalho espera contribuir com a área de ensino de programação sob quatro perspectivas complementares: teórica, técnica, empírica e pedagógica.

\subsection*{Contribuição teórica}

Um Mapeamento Sistemático da Literatura sobre o uso de logs no ensino de programação, com ênfase na produção nacional. Diferentemente de revisões anteriores, este mapeamento explicita não apenas as aplicações e ferramentas existentes, mas também as lacunas metodológicas e éticas ainda não endereçadas, oferecendo um panorama crítico que pode orientar futuras pesquisas na área.

\subsection*{Contribuição técnica}

O desenvolvimento e a disponibilização pública do plugin LOGADO para o Visual Studio Code, sob licença aberta. A contribuição não se limita ao código, mas inclui a documentação das decisões de projeto, dos eventos capturados e das justificativas pedagógicas para cada escolha, permitindo que outros pesquisadores adaptem e repliquem a solução em seus próprios contextos.

\subsection*{Contribuição empírica}

A disponibilização de um dataset com logs reais de programação, coletados em condições autênticas de ensino no Instituto Federal do Amazonas. O dataset é anonimizado e estruturado de forma a permitir análises replicáveis, respeitando a LGPD e os princípios da ciência aberta. Acredita-se que esse conjunto de dados possa subsidiar novos estudos sobre padrões de erro, estratégias de resolução e correlações entre comportamento de codificação e desempenho acadêmico.

\subsection*{Contribuição pedagógica}

Como contribuição adicional, este trabalho oferece uma reflexão sobre os limites da análise de logs no contexto educacional, discutindo o que a telemetria pode (e não pode) revelar sobre o processo de aprendizagem. Essa discussão pode ser útil para educadores que consideram adotar abordagens semelhantes, ajudando-os a calibrar expectativas e evitar interpretações equivocadas dos dados.

\section{Estrutura do trabalho}
\label{sec:estrutura-trabalho}

Este capítulo apresentou o contexto, a motivação e a justificativa do trabalho, suas questões de pesquisa, objetivos, contribuições esperadas e a estrutura da dissertação.

O Capítulo 2 expõe a fundamentação teórica, na qual são definidos os termos e conceitos fundamentais para o entendimento da proposta e discutidos os trabalhos correlatos.

O Capítulo 3 descreve a metodologia de pesquisa adotada, incluindo a abordagem metodológica (Mapeamento Sistemático, Survey Exploratório e Design Science Research), o plano de análise de dados e o percurso de execução da pesquisa.

O Capítulo 4 apresenta os resultados obtidos a partir da coleta e análise dos logs, com a caracterização da amostra, a estrutura do dataset gerado e as discussões sobre os padrões identificados.

O Capítulo 5 traz as considerações finais, retomando as contribuições do trabalho, suas limitações e apontando direções para pesquisas futuras.

Ao final, são apresentados os apêndices, que contêm:
\begin{itemize}
    \item o protocolo completo do Mapeamento Sistemático da Literatura (Apêndice A);
    \item o instrumento de coleta e os resultados do survey exploratório (Apêndice B);
    \item a estrutura e uma amostra do dataset gerado a partir dos logs coletados (Apêndice C).
\end{itemize}

O código-fonte do plugin LOGADO está disponível publicamente no repositório GitHub do author, em \url{https://github.com/gilmargn/logado-utfpr-ifam}, e será oportunamente registrado e incluído em versões futuras deste documento, em consonância com os princípios da ciência aberta.